\section*{Capítulo \ref{ch:b3:behavioural-ecology} --- Ecología del comportamiento}

\begin{solution}{exo:b3:behavioural-ecology:1}
Mecanismo: la danzante traslada el ángulo entre el alimento y el sol,
medido durante el vuelo, al ángulo respecto de la gravedad sobre el
panal, y codifica la distancia a partir del flujo óptico del viaje como
la duración del recorrido. Desarrollo: en gran medida innato ---las
abejas criadas sin ver danzas danzan correctamente, con pequeños
ajustes aprendidos. Función: recluta compañeras de nido hacia una
fuente provechosa y eleva el ingreso de la colonia. Historia: derivada
de movimientos de intención ritualizados de partida hacia el alimento;
las abejas enanas siguen danzando sobre una superficie horizontal
apuntando directamente a la fuente, la forma ancestral.
\end{solution}

\begin{solution}{exo:b3:behavioural-ecology:2}
Una \omterm{def:b3:behavioural-ecology:tinbergen}{pauta fija de acción} es una secuencia estereotipada que
desencadena una señal simple, el \omterm{def:b3:behavioural-ecology:tinbergen}{estímulo señal}, y que se completa una
vez comenzada: el rodar del huevo del ganso, desencadenado por un objeto
con forma de huevo; el picoteo del pollo de gaviota, desencadenado por
una mancha roja sobre una forma parecida a un pico. El palo con tres
bandas rojas, picoteado más que una cabeza de verdad, mostró que el
mecanismo desencadenante es un filtro afinado a un rasgo ---contraste
rojo, alargamiento--- y no un reconocimiento del progenitor, y que
exagerar el rasgo exagera la respuesta.
\end{solution}

\begin{solution}{exo:b3:behavioural-ecology:3}
Abandónese un rodal cuando el ritmo de ganancia en él haya caído al
ritmo medio del hábitat en su conjunto, viaje incluido. Duplicar el
tiempo de viaje baja esa media, de modo que el forrajeador se queda más
tiempo en todos los rodales, agota cada uno más a fondo y termina con
un ritmo global menor.
\end{solution}

\begin{solution}{exo:b3:behavioural-ecology:4}
Una EEE es una estrategia que, una vez que la juega casi todo el mundo,
ninguna alternativa rara puede superar. Todo Palomas no lo es: un solo
Halcón gana todos los enfrentamientos y obtiene $V$ frente al $V/2$ de
las Palomas. En la mezcla, cada estrategia obtiene $V(1 - V/C)/2$, menos
que el $V/2$ que daría todo Palomas, porque una fracción de los
enfrentamientos son peleas que hieren; nadie puede mejorarlo por su
cuenta, de modo que todos cargan con la pérdida.
\end{solution}

\begin{solution}{exo:b3:behavioural-ecology:5}
$V = 6$, $C = 10$: $p^{*} = 0.6$; pago $6(1 - 0.6)/2 = 1.2$ a cada una,
frente a $3$ de todo Palomas. $V = 12 > C$: Halcón es la EEE pura, con
pago $(12 - 10)/2 = 1$, frente a $6$ de todo Palomas ---el
enfrentamiento destruye cinco sextos del valor.
\end{solution}

\begin{solution}{exo:b3:behavioural-ecology:6}
$\tau = T_{0}$: $t^{*} = 1.15\times 2 = \qty{2.3}{min}$, fracción $1 -
\mathrm{e}^{-1.15} = 0.68$, $R = 0.68A/4.3 = 0.16A$ por minuto. $\tau =
4T_{0}$: $t^{*} = 2.0\times 2 = \qty{4}{min}$ (más exactamente $3.9$),
fracción $0.86$, $R = 0.86A/12 = 0.072A$ por minuto.
\end{solution}

\begin{solution}{exo:b3:behavioural-ecology:7}
Medio hermanos, $1/4$; abuelo y nieto, $1/4$; primos hermanos, $1/8$.
Una obrera haplodiploide y su hermano: él es haploide y solo tiene
genes de la madre de ambos; de los genes de la obrera, la mitad vinieron
de la madre y la mitad de esos son las copias que él recibió, de modo
que $r = 1/4$ desde el punto de vista de ella (desde el de él es $1/2$:
el parentesco no es simétrico cuando la ploidía difiere). Reina e hijo:
todos los genes de él son de ella, pero solo la mitad de los de ella
están en él: $r = 1/2$ por parte de ella y $1$ por parte de él.
\end{solution}

\begin{solution}{exo:b3:behavioural-ecology:8}
$n\,r\,b > c$ con $b = 0.01$ y $c = 0.02$: $n\,r > 2$. Hermanos
completos, con $r
= 1/2$: $n \ge 5$. Sobrinos, con $r = 1/4$: $n \ge 9$. Sin parentesco:
nunca ---y por eso los machos que se dispersan no llaman.
\end{solution}

\begin{solution}{exo:b3:behavioural-ecology:9}
Los apareamientos se duplican por cada \qty{25}{cm} y la supervivencia
cae una quinta parte: $L = 25$: $0.5\times 1.2 = 0.6$; $50$: $1$; $75$:
$2\times 0.8 = 1.6$; $100$: $4\times 0.6 = 2.4$; $125$:
$8\times 0.4 = 3.2$; $150$: $16\times 0.2 =
3.2$. En estos términos el producto sigue subiendo mucho más allá de
los \qty{50}{cm} naturales, de modo que al modelo le faltan costes: la
penalización de supervivencia de una cola larga en vuelo es más que
lineal, la respuesta femenina se satura, criar las plumas cuesta
energía y la duplicación de nidos de Andersson se midió a lo largo de
una temporada y no de una vida. La cola natural está allí donde muerden
los costes reales.
\end{solution}

\begin{solution}{exo:b3:behavioural-ecology:10}
Represalia: frente a Halcón $(V - C)/2$, frente a Paloma $V/2$ y frente
a Represalia $V/2$. En una población de solo Represalias, un Halcón raro
se encuentra con Represalias, que responden peleando, y obtiene
$(V - C)/2 < V/2$: no puede invadir cuando $V < C$. Una Paloma rara
obtiene $V/2$, exactamente el pago de las Represalias, de modo que las
Palomas son neutras y pueden colarse por deriva; una vez comunes, dejan
entrar a los Halcones (un Halcón obtiene $V$ frente a una Paloma), y el
juego de tres estrategias no tiene una EEE simple ---un primer ejemplo
de cómo añadir una opción puede desestabilizar una mezcla estable.
\end{solution}

\begin{solution}{exo:b3:behavioural-ecology:11}
Si todos persistieran exactamente $t_{0}$, un mutante que persistiese
$t_{0} +
\varepsilon$ ganaría todos los enfrentamientos por un coste adicional
despreciable, y a una población de tales mutantes la vencería
$t_{0} + 2\varepsilon$, y así sucesivamente: ningún tiempo fijo es
estable, de modo que la EEE ha de ser uno aleatorio. Con una
distribución exponencial, la probabilidad de persistir otro rato $s$,
dada una persistencia previa de $t$, es $\mathrm{e}^{-s/m}$ sea cual sea
$t$ ---la propiedad de falta de memoria---, de modo que un oponente que
sigue exhibiéndose no delata nada sobre cuánto más va a durar, y
ninguna estrategia puede aprovechar el tiempo transcurrido.
\end{solution}

\begin{solution}{exo:b3:behavioural-ecology:12}
Las obreras valoran a una hermana en $3/4$ y a un hermano en $1/4$, de
modo que bajo el control de las obreras la colonia debería invertir tres
veces más en hembras que en machos, una proporción de $3{:}1$ por
inversión; la reina, igual de emparentada con ambos, prefiere $1{:}1$.
Trivers y Hare (1976) hallaron proporciones de inversión cercanas a
$3{:}1$ en hormigas con una sola reina apareada una sola vez: ganan las
obreras. El apareamiento múltiple baja el parentesco medio de una obrera
con sus hermanas hacia $1/4$ o $1/2$ y empuja el óptimo de las obreras
hacia $1{:}1$; las colonias de especies con reinas apareadas muchas
veces invierten, en efecto, de forma más pareja, como predice el
razonamiento.
\end{solution}

\begin{solution}{pb:b3:behavioural-ecology:1}
\textbf{1.} $R(t) = A(1 - \mathrm{e}^{-t/T_{0}})/(\tau + t)$; márchese
cuando $g'(t) = R(t)$: $(A/T_{0})\mathrm{e}^{-t/T_{0}} = A(1 - \mathrm{e}^{-t/T_{0}})/
(\tau + t)$.
\textbf{2.} Multiplíquese por $(\tau + t)\mathrm{e}^{t/T_{0}}/A$: $(\tau + t)/T_{0}
= \mathrm{e}^{t/T_{0}} - 1$. Con $\tau = T_{0}$ y $x = t/T_{0}$: $2 + x =
\mathrm{e}^{x}$; $x = 1.15$ da $3.15$ frente a $3.16$.
\textbf{3.} $\tau = 4T_{0}$: $5 + x = \mathrm{e}^{x}$; $x = 1.94$ da
$6.94$ frente a $6.96$; $t^{*} \approx 1.94\,T_{0}$, digamos $2.0$.
\textbf{4.} Denso: $t^{*} = \qty{3.5}{min}$, $60(1 - \mathrm{e}^{-1.15}) =
\qty{41}{kJ}$, \qty{68}{\%}, $R = 41/6.5 = \qty{6.4}{kJ/min}$. Disperso:
$t^{*} = \qty{5.8}{min}$, \qty{51}{kJ}, \qty{86}{\%}, $R = 51/17.8 =
\qty{2.9}{kJ/min}$.
\textbf{5.} Cinco minutos: $g(5) = 60(1 - \mathrm{e}^{-5/3}) = \qty{49}{kJ}$.
Denso: $49/8 = \qty{6.1}{kJ/min}$ (el \qty{96}{\%} del óptimo);
disperso: $49/17 = \qty{2.9}{kJ/min}$ (el \qty{99}{\%}). Una regla burda
pierde un pequeño porcentaje ---bastante buena para que la selección se
haya conformado con ella.
\textbf{6.} El intervalo entre capturas se alarga a medida que el rodal
se agota, de modo que un tiempo de espera fijo es un ritmo marginal de
captura fijo ---el criterio del teorema, que es el mismo en todos los
rodales. Falla porque el umbral correcto depende del hábitat: con
\qty{20}{s} fijos el animal se marcha demasiado pronto donde el viaje
es largo y demasiado tarde donde es corto, salvo que el propio tiempo
de espera se ajuste a la experiencia reciente ---cosa que los
forrajeadores hacen.
\textbf{7.} Halcón contra Halcón, $(10 - 30)/2 = -10$; Halcón contra
Paloma, $10$; Paloma contra Halcón, $0$; Paloma contra Paloma, $5$;
$p^{*} = V/C = 1/3$.
\textbf{8.} $W_{H} = 10 - 20p$, $W_{D} = 5(1 - p)$. $p = 0.1$: $8$ frente
a $4.5$, se extienden los Halcones. $p = 1/3$: $3.33$ cada una,
equilibrio. $p = 0.6$: $-2$ frente a $2$, se extienden las Palomas.
\textbf{9.} $3.33$ en la EEE frente a $5$ de todo Palomas: un tercio del
valor de cada enfrentamiento se pierde por la posibilidad de pelear.
\textbf{10.} $V = 40 > C$: Halcón es la EEE pura, y todo enfrentamiento
es una pelea. Los años duros deberían traer más peleas y más heridas
---y así ocurre.
\textbf{11.} Un Halcón que se encuentra con un Burgués lo halla
residente la mitad de las veces (una pelea, $(V - C)/2$) e intruso la
otra mitad (un paseo, $V$): media $(3V - C)/4$. El Burgués obtiene $V/2$
entre los suyos. El Halcón invade solo si $(3V - C)/4 > V/2$, es decir,
$V > C$. El ``gana el residente'' de la mariposa es Burgués: una
convención barata, estable porque pelear cuesta más de lo que vale un
claro de sol ---y cuando ambos se creen residentes, ambos juegan a
Halcón.
\textbf{12.} Lo mejor que se puede hacer depende de lo que hagan todos
los demás, de modo que la población se asienta allí donde las
estrategias pagan igual y no en el óptimo de ningún individuo.
\textbf{13.} $n\bar r\,b > c$: $n\bar r > c/b = 3$.
\textbf{14.} Ocho hermanos completos: $n\bar r = 4 > 3$, la guardia
compensa. Ocho desconocidos: $0$, nunca. Ocho medio hermanos: $2 < 3$,
no.
\textbf{15.} $c = 0.01$: umbral $n\bar r > 1$ ---ahora bastan los medio
hermanos e incluso unos pocos hermanos completos. El individuo con el
coste menor es aquel para el que se cumple la regla, de modo que los
bien alimentados montan guardia mientras los hambrientos escarban.
\textbf{16.} $k$ hermanas valen $\tfrac{3}{4}k$ y $k$ hijas
$\tfrac{1}{2}k$: prefiéranse las hermanas, en proporción $3{:}2$.
\textbf{17.} Una hermana al azar es hermana completa ($3/4$) con
probabilidad $1/3$ y media hermana ($1/4$) con probabilidad $2/3$: $\bar r = 0.25 +
0.17 = 0.42 < 0.5$. La preferencia se invierte: las hijas superan ahora
a las hermanas, de modo que el parentesco por sí solo no puede explicar
las obreras estériles en las especies poliándricas ---hay que añadir la
ecología y la vigilancia a escala de colonia.
\textbf{18.} La regla solo exige que $r$ sea mayor, en promedio, entre
los ayudados que en la población; sirve cualquier señal correlacionada
con el parentesco ---quedarse cerca de la madriguera natal, ayudar a
aquellos con los que se ha crecido. La regla de una ardilla terrestre
puede ser ``llama cuando estés cerca de casa'', que selecciona parientes
sin reconocer a ninguno.
\textbf{19.} $W = (L/20)(1 - L^{2}/10^{4})$: $L = 20$: $0.96$; $40$:
$1.68$; $60$: $1.92$; $80$: $1.44$.
\textbf{20.} $W' = \tfrac{1}{20}(1 - 3L^{2}/10^{4}) = 0$: $L^{*} = 100/\sqrt{3}
= \qty{58}{cm}$, $W = 2.89\times\tfrac{2}{3} = 1.92$.
\textbf{21.} $L^{*} = 70/\sqrt{3} = \qty{40}{cm}$: más corta. El óptimo
se sitúa allí donde la ganancia marginal en apareamientos iguala a la
pérdida marginal en supervivencia, de modo que un coste de depredación
mayor encoge la cola, y las poblaciones sometidas a depredaciones
distintas deberían diferir en el ornamento ---como ocurre en los
guppys y en los xifos.
\textbf{22.} Con ambas. La \omterm{def:b3:behavioural-ecology:sexual-selection}{selección desbocada} predice que la
preferencia se adelanta al rasgo, que se detiene en el óptimo de
supervivencia; el hándicap predice que las hembras prefieren las más
largas porque solo un macho en forma la lleva, de nuevo más allá de lo
que la mayoría puede criar. El experimento muestra una preferencia más
allá del óptimo y no decide entre las dos explicaciones.
\textbf{23.} Si la cola fuera gratis, todos los machos podrían criar la
más larga, no diría nada sobre su calidad y una preferencia por ella
quedaría explotada y decaería. Un coste que pesa más sobre un macho en
mala condición hace la cola larga asequible solo para los que están en
forma, y la preferencia selecciona entonces eficacia: el coste es lo que
vuelve honesta la señal.
\textbf{24.} Con ocho veces el rendimiento por apareamiento, la eficacia
de un macho sube de forma empinada con los apareamientos y la de una
hembra apenas sube; los machos quedan seleccionados para competir por
apareamientos ---los enfrentamientos Halcón--Paloma y sus armas--- y las
hembras, cuyos pocos descendientes cuentan cada uno, quedan
seleccionadas para elegir, lo que hace rentables los ornamentos de los
machos. Una sola asimetría en la ganancia por apareamiento produce a la
vez las peleas y las galas.
\textbf{25.} $t^{*} \approx \qty{3.5}{min}$ en el hábitat denso; $p^{*}
= 1/3$; la guardia compensa cuando $n\bar r > 3$.
\end{solution}
